\documentclass{article}

% if you need to pass options to natbib, use, e.g.:
%     \PassOptionsToPackage{numbers, compress}{natbib}
% before loading neurips_2026

%\usepackage{neurips_2026}

% 正确传递 natbib 选项，不冲突
\PassOptionsToPackage{numbers, sort&compress}{natbib}
\usepackage[dandb, final]{neurips_2026}

\usepackage{graphicx}
\usepackage{float}
\usepackage{listings}
\usepackage{xcolor}

\lstset{
  basicstyle=\ttfamily\small,
  breaklines=true,
  frame=single,
  columns=fullflexible,
  keepspaces=true,
  showstringspaces=false
}
 
\usepackage[utf8]{inputenc}
\usepackage[T1]{fontenc}
\usepackage{hyperref}
\usepackage{url}
\usepackage{booktabs}
\usepackage{amsfonts}
\usepackage{amssymb}
\usepackage{amsmath}
\usepackage{amsthm}
\usepackage{nicefrac}
\usepackage{microtype}
\usepackage{xcolor}
\usepackage{array}
\usepackage{algorithm}
\usepackage{algpseudocode}
\usepackage{graphicx}
\newcommand{\gx}[1]{\textcolor{cyan}{#1}}
\newcommand{\wang}[1]{\textcolor{blue}{#1}}
\newcommand{\wei}[1]{\textcolor{red}{#1}}

\title{HSPO: Hierarchical Subgoal-conditioned Process Optimization for Long-Horizon LLM Agents}

\author{
Zhuoting Wang$^{1,2}$, Letong Meng$^{1,2}$, 
\textbf{Wei Du}$^2$\textbf{,} \textbf{Dong Wu}$^3$\textbf{,} \textbf{Dong Zhang}$^3$\textbf{,} \textbf{Jun Wang}$^1$\textbf{,} \textbf{Guoxian Yu}$^{1,}$\thanks{Corresponding author} \\
$^1$School of Software, Shandong University, Jinan, China \\
$^2$SDU-NTU Centre for AI Research, Shandong University, Jinan, China \\
$^3$Inspur Data Technology Co, Ltd, Jinan, China\\
\texttt{\{zhtwang, ltmeng\}@mail.sdu.edu.cn},  \texttt{\{zhangdong, wudongbj\}@inspur.com}, \\ \texttt{\{duwei, kingjun, gxyu\}@sdu.edu.cn}
}

\begin{document}

\maketitle
\begin{abstract}
Reinforcement learning (RL) for long-horizon LLM agents typically learns from sparse, trajectory-level returns, which conflate strategic subgoal selection with low-level action execution. When a chosen subgoal is globally suboptimal, all actions in its segment can be penalized, including locally correct behaviors that should remain reusable under a better plan. We call this failure mode as \textbf{intra-subgoal credit collapse}, and propose a hierarchical RL framework \textbf{Hierarchical Subgoal-conditioned Process Optimization} (\textbf{HSPO}) that mitigates this collapse via three coordinated mechanisms. (i) An online Plan-Execute interface factorize each step into termination, subgoal selection, and action execution, so credit is assigned at appropriate decision level. (ii)  \textbf{Anchored Branch Grouping (ABG)} contrasts candidate actions under identical state-subgoal anchors and scores them with a state--grounded milestone scorer, isolating local execution progress from delayed global outcomes. (iii) Token-level credit routing applies termination, high-level and low-level losses only to their corresponding output spans, and adds subgoal-distribution regularization to prevent planner drift during executor updates. This objective discourages globally harmful subgoals while preserving executor behaviors that correctly implement them.
\end{abstract}


%\wei{摘要参考版本：Reinforcement learning for long-horizon LLM agents relies on flat trajectory‑level returns, which conflate subgoal selection and primitive action execution.(背景，一句话)  This leads to intra‑subgoal credit collapse, where when a subgoal is globally suboptimal, every action within that segment is penalized, including locally correct executions that should remain reusable under better subgoal choices. (问题，用你定义的这个信用坍塌)  We propose Hierarchical Subgoal-conditioned Process Optimization (HSPO), a hierarchical RL framework that resolves this collapse through three mechanisms. First, an explicit Plan‑Execute interface (<switch>/<subgoal>/<action>) separates subgoal selection from action execution at every step, ensuring credit can be routed to the correct level. Second, a subgoal‑grounded local credit mechanism constructs action comparison groups under fixed task–state–subgoal anchors and scores each action by deterministic one‑step progress toward the active subgoal, so that locally correct actions are no longer penalized by a globally bad subgoal. Third, token‑level credit routing applies high‑level, low‑level, and termination losses only to their corresponding spans, paired with a KL regularizer that prevents planner drift during executor updates, thereby decoupling the two levels of optimization. 方法，三个内容不用展开，概述解决了什么问题)The result is that the planner can be discouraged from choosing a harmful subgoal, while the executor simultaneously preserves actions that correctly implement it.(结果，目前是定性的概括，后面加实验结果)}




\section{Introduction}

Large language models (LLMs) are increasingly deployed as interactive agents that observe their environments, plan over multiple steps, and execute actions in dynamic tasks~\cite{yao2022react,shinn2023reflexion,wang2024voyager}. Reinforcement learning (RL) provides a natural framework for improving such agents from interaction feedback. However, in long-horizon settings rewards are often sparse or delayed, making it hard to assign credit across the heterogeneous decisions within a trajectory.

%\wei{第一段很好，介绍背景和挑战，不用改}

A common way to address sparse feedback is to optimize the policy with trajectory-level or segment-level returns. However, this approach breaks down when the agent behaves hierarchically, i.e., when it first chooses an intermediate subgoal and then executes a sequence of primitive actions to achieve that subgoal before switching to the next. Consider the task \emph{clean a cup and place it in a cabinet}. If the agent chooses the subgoal \emph{cool the cup}, the subgoal is globally inappropriate for the task. Yet, conditioned on this active subgoal, the action \emph{put the cup in the fridge} is a reasonable local execution. A trajectory- or segment-level return penalizes the whole segment uniformly, thereby mixing two different questions: whether the subgoal should have been selected, and whether the executor correctly carried it out. As a result, the agent may learn to suppress useful low-level behaviors simply because they were performed under a bad high-level plan.


%已改
%\wei{第二段举例说明要解决的问题。这一段也可以，但是没定义hierarchical agent behavior就开始举例子，同时最后一句表达的不够准确。参考版本：A simple and widely used approach is to optimize the policy with trajectory‑level or segment‑level returns. However, this approach breaks down when the agent behaves hierarchically, that is, when it first chooses a subgoal and then executes a sequence of primitive actions to achieve that subgoal before moving to the next. Consider the task \emph{clean a cup and place it in a cabinet.} If the agent chooses the subgoal \emph{cool the cup}, that subgoal is globally wrong. Yet, given that active subgoal, the action \emph{put the cup in the fridge} is locally correct. 最后一句调整得更易懂一些,因为审稿人通常是大同行：A trajectory return penalizes the whole segment equally, which mixes together two different questions: “Is this subgoal a good choice?” and “Given this subgoal, did I execute the action properly?” The agent thus learns to avoid good low‑level actions simply because they were performed under a bad high‑level plan. }




We call this failure mode \textbf{intra-subgoal credit collapse}. It occurs when the global utility of a selected subgoal is inherited by every primitive action within its execution segment. Consequently, the executor may unlearn reusable skills simply because they were invoked under an incorrect global plan. A robust hierarchical agent therefore requires not only a decomposition of decisions into subgoal selection and action execution, but also a credit assignment mechanism that evaluates these two decisions independently.

\begin{figure*}[t]
\centering
\includegraphics[width=\textwidth]{figures/fig1-compare.png}
\caption{
\textbf{Illustration of intra-subgoal credit collapse and HSPO's credit separation.}
(a) Traditional segment-level RL propagates a negative return to all actions inside a failed segment. When the agent selects a globally inappropriate subgoal such as \emph{cool the cup}, even a locally correct action under that subgoal, e.g., putting the cup in the fridge, is penalized. This causes the executor to forget reusable low-level skills.
(b) HSPO separates strategic subgoal selection from subgoal-conditioned action execution. The planner is penalized for choosing a harmful subgoal, while the executor can still receive positive low-level credit for actions that correctly advance the active subgoal. Anchored Branch Grouping (ABG), milestone scoring, and token-level credit routing jointly prevent globally bad subgoals from collapsing credit onto locally correct actions.
}
\label{fig:credit_collapse}
\end{figure*}

%已改
%\wei{第三段，定义问题intra-subgoal credit collapse。这段前面写的也不错，最后一句有些凝练了，你现在的写作风格，是只能给熟悉你这个方案或者小同行的人看懂，写文章的目标是所有大同行看到都一下子能明白你的动机，尤其引言部分。然后又觉得你想到的创新点，是他们没想到的，不容易想的，这种效果是最好。We call this failure mode \textbf{intra-subgoal credit collapse: the global utility of a subgoal is inherited by every primitive action inside its execution segment. Consequently, the executor may unlearn reusable skills simply because they were invoked under a wrong global plan. A robust hierarchical agent therefore requires not only a hierarchical decomposition of decisions, but also a credit assignment mechanism that judges subgoal selection and action execution independently.}

Recent methods have made important progress toward improving credit assignment for long-horizon LLM agents. GRPO~\cite{shao2024deepseekmath} estimates relative advantages from groups of sampled trajectories and avoids learning an explicit value critic. However, because its comparison groups are defined at the trajectory level, low-level actions still inherit outcome-level credit from final task success or failure. HiPER~\cite{wang2026subgoal} introduces an online Plan-Execute interface and dynamically segments trajectories into subgoal-conditioned executions. Yet, if actions within each segment are still optimized with segment-level or trajectory-level returns, locally correct executions can still be penalized when their parent subgoal is strategically wrong. GiGPO~\cite{NEURIPS2025_420c9f77} constructs step-level comparison groups from repeated environment states across rollouts, providing finer-grained feedback than trajectory-level objectives. However, its anchors are defined mainly by state, without explicitly conditioning on the active subgoal; the same state-action pair may therefore receive the same credit signal even when it serves different intents. HGPO~\cite{wang2025spa} further improves grouped policy optimization by addressing historical-context inconsistency, but it still does not define a subgoal-grounded process signal that directly judges whether an action advances the currently active subgoal.

%\wei{第四段，先前工作和缺陷。这段写作问题就比较明显了，讲完别人的进展，戛然而止了，中间少了一段关键的解释，为什么他们不好，分析的比较浅，这里我等下看看补充一段。另外，这一段也介绍一下GRPO吧，还有HiPER这种介绍太短了，然后这段的最后一句的总结，看起来不太好直接和前面几个方法存在的问题很直观的呼应，得读半天才对得上。冒号这种表达，除了你定义credit collapse那里可以用一下，尽量不要使用，最好每个方法存在的缺陷，直接接着这个方法讲一下就行，下面我新增加的这一段再总结性的表达。}

More generally, existing RL methods for LLM agents either redistribute final outcomes to earlier decisions or construct comparison groups using trajectory- or state-level anchors. Both strategies remain insufficient for hierarchical behavior. When low-level credit depends on final task success, a locally correct action under a globally harmful subgoal can still receive negative credit. When grouping ignores the active subgoal, actions are compared without knowing which intent they are supposed to execute. In both cases, the executor may unlearn reusable skills simply because those skills were invoked under the wrong high-level plan. This is precisely the intra-subgoal credit collapse that our method is designed to resolve.


%\wei{这里补充第五段，承上启下。More generally, prior hierarchical RL for LLM agents either derives low‑level credit from segment‑level returns, which propagates the global subgoal utility to every action, or relies on outcome‑based reward redistribution that does not break dependency on final success. Group‑based estimators such as GRPO \cite{} and GiGPO eliminate the value critic but still compare actions using trajectory outcomes or state‑level anchors, without conditioning on the active subgoal. Thus, a locally correct action that advances a temporarily active but globally harmful subgoal still receives negative credit, and the executor unlearns a reusable skill. This is precisely the intra‑subgoal credit collapse that our method is designed to resolve.}



To resolve this collapse, we propose \textbf{Hierarchical Subgoal-conditioned Process Optimization} (\textbf{HSPO}), a hierarchical RL framework that separates subgoal selection from subgoal-conditioned action execution. HSPO uses an online Plan-Execute interface in which the policy emits
\[
\texttt{<switch>}q_t\texttt{</switch>}
\texttt{<subgoal>}g_t\texttt{</subgoal>}
\texttt{<action>}a_t\texttt{</action>}
\]
at every environment step. The switch span induces subgoal segmentation online; the subgoal span represents high-level intent; the action span executes the current subgoal. 

The core of HSPO is \textbf{Anchored Branch Grouping} (ABG), a low-level group-construction mechanism that compares actions only under the same task--state--subgoal anchor. After collecting complete task rollouts, ABG scans all low-level decision points and groups transitions that share the same task signature, state signature, and active subgoal instance. Naturally repeated anchors are reused directly, while sparse groups are selectively supplemented with one-step branched candidate actions. Each candidate transition is then evaluated by a deterministic, state-grounded milestone scorer that measures one-step progress toward the active subgoal, rather than delayed progress toward the final task outcome.

HSPO routes different credit signals to different spans: high-level macro advantages update only \texttt{<subgoal>} tokens, low-level ABG advantages update only \texttt{<action>} tokens, and termination supervision updates only \texttt{<switch>} tokens.

As a result, HSPO separates the source of failure from the level of behavior being optimized. When the overall task fails because the agent follows an inappropriate subgoal plan, the negative credit should primarily affect the high-level planner that selected those subgoals. The low-level executor, however, should be optimized according to how well it completes the currently active subgoal, rather than whether that subgoal ultimately contributes to task success. In this way, high-level learning is driven by the quality of the planned subgoal sequence, while low-level learning is driven by subgoal-conditioned execution progress. HSPO therefore discourages globally harmful subgoal plans without erasing reusable low-level skills that correctly implement a given subgoal.

%【已加】
%\wei{第六段，引出我们的方法，并详细展开介绍了，这里就是摘要的plus版本，越细致越好，然后最好的情况，可以配一个图，和前面举的那个例子一起，左边是clean a cup and place it in a cabinet分层动作演示，右边是我们的方法的几个创新点和解决后的大概预期效果。第六段把你这三段并到一起了，写的挺好，先这样并一下，分了123, 后续可以再微调一下表述：To resolve this collapse, we propose Hierarchical Subgoal-conditioned Process Optimization (\textbf{HSPO}). (1) HSPO uses an online Plan-Execute interface in which the policy emits \[\texttt{<switch>}q_t\texttt{</switch>}\texttt{<subgoal>}g_t\texttt{</subgoal>}\texttt{<action>}a_t\texttt{</action>}\] at every environment step. The switch span induces subgoal segmentation online; the subgoal span represents high-level intent; the action span executes the current subgoal. (2) The core of HSPO is \textbf{Anchored Branch Grouping} (ABG), a low-level group-construction mechanism. After collecting complete task rollouts, ABG scans all low-level decision points and groups transitions with the same task signature, state signature, and active subgoal instance. Naturally repeated anchors are reused directly, while sparse groups are selectively supplemented with one-step branched candidate actions. Each candidate transition is scored by a deterministic state-grounded milestone scorer that measures progress toward the active subgoal. (3) HSPO routes different credit signals to different spans: high-level macro advantages update only \texttt{<subgoal>} tokens, low-level ABG advantages update only \texttt{<action>} tokens, and termination supervision updates only \texttt{<switch>} tokens. Thus, a planner can be penalized for selecting a globally harmful subgoal, while the executor can still be rewarded for actions that correctly implement that subgoal.

%【已加】
%\wei{这里增加一个第七段，总结一下我们的方法的最终贡献。As a result, the planner can be discouraged from choosing a globally harmful subgoal, while the executor simultaneously preserves actions that correctly implement that subgoal. This resolves intra‑subgoal credit collapse because a locally correct action under a bad subgoal receives positive low‑level advantage, since it advances the subgoal’s internal progress, while the subgoal itself receives negative high‑level advantage, since it leads to overall task failure. The two signals are routed to disjoint token spans and do not interfere.}

Our contributions are summarized as follows:
\begin{itemize}
    \item We identify and formalize \textbf{intra-subgoal credit collapse} as a distinct failure mode in long-horizon RL for LLM agents, where low-level reusable skills are incorrectly penalized due to globally harmful subgoal choices.
    \item We propose \textbf{Anchored Branch Grouping} (ABG), which constructs task--state--subgoal comparison groups and uses deterministic state-grounded milestone scorers to provide subgoal-conditioned one-step process feedback.
    \item We introduce a unified hierarchical RL framework, \textbf{HSPO}, that combines online Plan-Execute segmentation, high-level macro optimization, low-level ABG optimization, lagged switch supervision, and token-level credit routing to decouple strategic subgoal selection from tactical execution.
\end{itemize}

%\wei{这里通常是分三个点，是一个约定俗成的写法。我给整合一下，\item We identify and formalize \textbf{intra-subgoal credit collapse} as a distinct failure mode in long-horizon LLM-agent RL. \item \textbf{Anchored Branch Grouping} (ABG) with deterministic milestone scorers to provide subgoal‑grounded, one‑step process feedback. \item We introduce a \textbf{unified hierarchical RL framework} that decouples strategic subgoal selection from tactical execution via token‑level credit routing.


\section{Related Work} \wei{Note: 我看最后引用的参考文献较少，这块可以让乐童帮着调研一些，多补充一些相关的工作进来}
\label{sec:relwork}

%\meng{先在注释里放一下想补充的文献:LLM agents部分：1.AgentBoard: An analytical evaluation board of multi-turn llm agents，证明随着交互轮数增加，大模型智能体的成功率会断崖式下跌，产生严重的级联错误。（第13篇，在实验部分引用过）2.TravelPlanner: A Benchmark for Real-World Planning with Language Models ，ICML 2024，证明在面对多约束、长周期的真实任务时，即使是 GPT-4 这种最强模型，如果只靠 Prompt 里的思维链（CoT）来规划，也是极其脆弱的  3.SwiftSage: A Generative Agent with Fast and Slow Thinking for Complex Interactive Tasks  NIPS 2023 尝试“快慢系统（高层规划+底层执行）分离”的先驱，为了肯定“结构化拆分”的正确性 4.AdaPlanner: Adaptive Planning from Feedback with Language Models NIPS 2023，证明前人已经意识到要根据环境反馈来动态调整 Plan，但依然是在推理阶段做 prompt 层面的调整，没有深入到强化学习的 Token 级别。5.AgentBench: Evaluating LLMs as Agents  ICLR 2024 证明“大模型作为智能体在多轮任务中表现极其拉垮”的一座大山 6.Tree of Thoughts: Deliberate Problem Solving with Large Language Models NIPS 2023 证明“复杂的任务必须用树状/非线性结构来拆解规划，不能顺着一直写”，支持扁平的自回归生成将高层意图和底层执行纠缠这一论断 7.MemGPT: Towards LLMs as Operating Systems  ICLR 2024 长视距任务之所以会产生“级联错误” ，不仅是因为规划做不好，还因为上下文太长导致模型遗忘

\textbf{LLM agents.}
Large language models (LLMs) are increasingly deployed as interactive decision-making agents for tool use, web navigation, embodied control, software engineering, and open-ended games~\cite{yao2022react,shinn2023reflexion,wang2024voyager}. Early agents typically relied on frozen foundation models steered via prompting, such as ReAct~\cite{yao2022react} and Reflexion~\cite{shinn2023reflexion}, often augmented with memory, retrieval, or external tools. While flexible, flat autoregressive generation entangles high-level intent with low-level execution in a single token stream~\cite{NEURIPS2023_271db992}, making long-horizon behavior prone to cascading errors and context-length bottlenecks~\cite{packer2024memgptllmsoperatingsystems}. This fragility is evidenced by severe performance degradation in recent multi-turn and long-horizon benchmarks under pure prompting~\cite{ma2024agentboard,ICLR2024_e9df36b2,pmlr-v235-xie24j}. Subsequent work has therefore explored structured agent interfaces—such as decoupling planning from execution or utilizing adaptive feedback~\cite{NEURIPS2023_4b0eea69,NEURIPS2023_b5c8c1c1}—alongside supervised fine-tuning~\cite{lightman2023let,wang2026subgoal} and reinforcement learning~\cite{schulman2017proximal,NEURIPS2025_420c9f77,ji2026tree} to improve multi-step decision making. HSPO follows this direction by exposing termination, subgoal selection, and action execution through an online Plan-Execute interface.

% \textbf{LLM agents.}
% Large language models (LLMs) are increasingly deployed as interactive decision-making agents for tool use, web navigation, embodied control, software engineering, and open-ended games~\cite{yao2022react,shinn2023reflexion,wang2024voyager}. Early agents typically relied on frozen foundation models steered via prompting, such as ReAct~\cite{yao2022react} and Reflexion~\cite{shinn2023reflexion}, often augmented with memory, retrieval, or external tools. While flexible, flat autoregressive generation entangles high-level intent with low-level execution in a single token stream, making long-horizon behavior prone to cascading errors. Subsequent work has therefore explored structured agent interfaces~\cite{xxx,yyy}, supervised fine-tuning~\cite{xxx,yyy}, and reinforcement learning~\cite{xxx,yyy} to improve multi-step decision making. HSPO follows this direction by exposing termination, subgoal selection, and action execution through an online Plan-Execute interface.
%\wei{相关工作这里写得还不错，但是我考虑这三个分类方式是不是有些太宽泛了，llm agent，rl for llm和hrl。这部分是全文最不重要的部分，可以先这样，最后再处理。}}


% \textbf{RL for LLM agents.}
% RL is a widely used paradigm for adapting LLM agents from environment feedback. Classical actor--critic methods such as PPO~\cite{schulman2017ppo} have been applied to embodied and web-based agents, but accurate value learning over high-dimensional, language-conditioned states remains difficult under sparse rewards. Critic-free, group-based objectives---including RLOO~\cite{ahmadian2024back}, GRPO~\cite{shao2024deepseekmath}, and DAPO~\cite{yu2025dapo}---alleviate this challenge by estimating relative advantages within sampled groups. More recent agentic RL methods further refine credit assignment in multi-turn settings: GiGPO~\cite{gigpo} forms step-level comparison groups from repeated environment states across rollouts, while HGPO~\cite{hgpo} addresses historical-context inconsistency in grouped policy optimization. Although these approaches yield denser, more stable training signals, they typically compare actions at the state or trajectory level without explicitly conditioning low-level credit on the active subgoal. In contrast, HSPO constructs action groups under shared task--state--subgoal anchors and scores one-step progress toward the current subgoal.

%这里gemini提出了一个新的段落名
\textbf{Process-based Optimization and Credit Assignment.} 
Reinforcement learning (RL) is a pivotal paradigm for adapting LLM agents through environment feedback, with foundational frameworks like InstructGPT~\cite{NEURIPS2022_b1efde53} establishing the standard for aligning models. However, as noted by Ramamurthy et al.~\cite{ramamurthy2023is}, applying RL to language-based tasks faces inherent complexities in reward density and policy optimization. Classical actor-critic methods, such as PPO~\cite{schulman2017proximal}, often struggle in long-horizon embodied or web-based environments because learning accurate value functions over high-dimensional states is exceptionally difficult under sparse, trajectory-level rewards~\cite{NEURIPS2025_420c9f77}.

To mitigate this sparsity, recent research has pivoted toward step-level evaluation and Process Reward Models (PRMs)~\cite{lightman2024lets}. For instance, LATS~\cite{pmlr-v235-zhou24r} integrates tree search with LM-powered value functions, while Math-Shepherd~\cite{wang-etal-2024-math} utilizes automated step-wise verification to bypass the need for human annotation. Other approaches, such as Motif~\cite{klissarov2024motif}, leverage AI feedback to generate intrinsic motivation for agent exploration. Nevertheless, these process-supervision methods often rely on LLM-as-a-judge mechanisms that are computationally expensive, prone to hallucinations, and not strictly grounded in the underlying physical environment state.

Concurrently, critic-free group-based objectives---including RLOO~\cite{ahmadian2024back}, GRPO~\cite{shao2024deepseekmath}, and DAPO~\cite{yu2025dapo}---alleviate the burden of value estimation by calculating relative advantages within sampled trajectories. Recent agentic RL methods like GiGPO~\cite{NEURIPS2025_420c9f77} and HGPO~\cite{wang2025spa} further improve credit assignment by constructing comparison groups from repeated environment states. However, these methods generally compare actions at the state or trajectory level without explicitly conditioning credit on the active subgoal. In contrast, HSPO introduces a deterministic, state-grounded milestone scorer and constructs action groups under shared task--state--subgoal anchors. This allows for objective scoring of one-step progress toward a specific intent, effectively preventing the \textit{intra-subgoal credit collapse} that plagues prior methods.


\textbf{Hierarchical reinforcement learning.}
Hierarchical RL addresses long-horizon decision making through temporal abstraction. Classical frameworks such as the options framework~\cite{sutton1999between} and feudal reinforcement learning~\cite{vezhnevets2017feudal} decompose behavior into high-level goal selection and low-level action execution. Recent LLM-agent methods adapt these ideas to natural-language interfaces. HiPER~\cite{wang2026subgoal} performs online subgoal segmentation during interaction, enabling hierarchical credit assignment over dynamically induced segments, while other subgoal-driven agents decouple macro-level planning from micro-level execution over natural-language goals or blueprints~\cite{ji2026tree}. Although such hierarchies improve long-horizon behavior, low-level actions can still inherit credit signals dominated by segment- or trajectory-level outcomes. HSPO targets this gap by separating strategic subgoal credit from tactical execution credit via ABG, deterministic state-grounded milestone scoring, and token-level credit routing, allowing the planner to avoid harmful subgoals while preserving executor behaviors that correctly implement the chosen subgoal.


\section{Method}
\label{sec:method}
HSPO disentangles two forms of credit that are conflated in flat, long-horizon agent RL: (i) whether a subgoal is strategically appropriate for the overall task, and (ii) whether a primitive action correctly executes the currently active subgoal. This section presents the complete method. We first define the online Plan-Execute interface, then introduce the state-grounded process scorer, Anchored Branch Grouping (ABG) for the low-level objective, and the high-level and termination objectives. Finally, we describe token-level credit routing and our staged training procedure. Figure~\ref{fig:hspo_overview} provides an overview of the full HSPO training pipeline.%\wei{这段写得很好}

\begin{figure*}[t]
\centering
\includegraphics[width=\textwidth]{figures/fig2-overview.png}
\caption{
\textbf{Overview of HSPO.}
HSPO optimizes long-horizon LLM agents through a hierarchical subgoal-conditioned process objective.
}
\label{fig:hspo_overview}
\end{figure*}


\subsection{Problem Formulation} %\wei{3.1很好，没发现问题}
\label{sec:problem_formulation}

We consider an interactive language-agent environment where a task instruction $x \sim \mathcal{D}$ is sampled from a task distribution. At each step $t$, the agent observes a state $s_t$, generates a structured response $y_t$, executes an environment action $a_t$, and transitions to $s_{t+1}$. Feedback is sparse: the main task reward $R_{\mathrm{task}}$ is typically observed only at the end of the episode.

HSPO adopts an online Plan-Execute interface that constrains each response to three fields:
\begin{equation}
\begin{aligned}
y_t
=
\langle q_t,g_t,a_t\rangle
=
&\ \texttt{<switch>}q_t\texttt{</switch>}
\texttt{<subgoal>}g_t\texttt{</subgoal>} \\
&\ \texttt{<action>}a_t\texttt{</action>}.
\end{aligned}
\label{eq:response}
\end{equation}
Here, $q_t \in \{\mathsf{Keep}, \mathsf{Switch}\}$ is a switch decision, $g_t$ is the active natural-language subgoal, and $a_t$ is the environment action. If $q_t=\mathsf{Keep}$, the agent continues executing the current subgoal; if $q_t=\mathsf{Switch}$, it initiates a new subgoal.

Let $0=b_0<b_1<\cdots<b_K\leq T$ denote the time indices at which the agent initiates new subgoals. These switch points induce a sequence of subgoal-conditioned segments. The $k$-th segment is
\begin{equation}
\tau_k
=
(s_{b_k},a_{b_k},s_{b_k+1},\ldots,a_{b_{k+1}-1},s_{b_{k+1}}),
\label{eq:segment}
\end{equation}
with active subgoal $g_k$. Thus, the trajectory is decomposed online into temporally extended execution segments indexed by subgoals.

Although a single autoregressive LLM $\pi_\theta$ emits all tokens in $y_t$, we conceptually separate its outputs into three decision channels for credit assignment:
\begin{equation}
\pi_\theta^T(q_t \mid x,s_t,g_{t-1},h_t),\quad
\pi_\theta^H(g_t \mid x,s_t,h_t),\quad
\pi_\theta^L(a_t \mid x,s_t,g_t,h_t),
\label{eq:factorization}
\end{equation}
where $h_t$ denotes a compact interaction history. The superscripts $T$, $H$, and $L$ correspond to switching, high-level subgoal selection, and low-level action execution, respectively. This factorization is used only to define losses and token-level credit routing during training; at inference time, HSPO remains a single online autoregressive agent.

%！！！加一个overview的图【图2-Architectural Overview】


\subsection{Plan-Execute Supervised Warm-up} %\wei{3.2也很好，没发现问题，你很擅长写比较细的方案}

HSPO begins with a lightweight supervised warm-up that teaches the model to reliably emit the structured Plan-Execute interface before RL. Successful demonstration trajectories are converted into step-level targets. When environment-provided subgoal annotations are available, we use them to define segment boundaries; otherwise, a constrained segmenter maps the trajectory to a finite canonical subgoal vocabulary $\Sigma$. We avoid unconstrained free-form segmentation, since a drifting subgoal vocabulary would make the process scorer and ABG anchors unreliable.

For each step $t\in[b_k,b_{k+1})$, the target response is
\begin{equation}
y_t =
\texttt{<switch>}q_t\texttt{</switch>}
\texttt{<subgoal>}g_k\texttt{</subgoal>}
\texttt{<action>}a_t\texttt{</action>},
\end{equation}
where $q_{b_k}=\mathsf{Switch}$ and $q_t=\mathsf{Keep}$ for $t>b_k$. We train with standard behavior cloning, i.e., next-token negative log-likelihood on demonstration responses:
\begin{equation}
\mathcal{L}_{\mathrm{SFT}}(\theta)
=
-\mathbb{E}_{(c,y)}
\sum_i
\log \pi_\theta(y_i\mid y_{<i},c),
\label{eq:sft_loss}
\end{equation}
where $c$ denotes the context. This SFT checkpoint stabilizes the Plan-Execute interface and is used to initialize all trainable Plan-Execute baselines.Full prompt templates for ALFWorld and WebShop are provided in Appendix~\ref{app:prompts}.


\subsection{State-Grounded Milestone Scorer} %\IrreversibleHarm
% 这里待会还再细化一下 OK 后面我一会再看，这个IrreversibleHarm有点怪，是自己起的名字吗，看看能不能换个词，或者解释一下。
The low-level executor should be evaluated by whether its action makes progress toward the active subgoal, rather than by whether that subgoal is globally useful. We therefore define a deterministic, state-grounded scorer:
\begin{equation}
\begin{aligned}
S(x,g,s,a,s')
=
&\ \Delta P_g(s,a,s')
+ \eta_{\mathrm{done}}\mathbf{1}[D_g(s')]  \\
&-\lambda_{\mathrm{inv}}\mathbf{1}[\mathrm{Invalid}(a,s)]
-\lambda_{\mathrm{hard}}\mathbf{1}[\mathrm{IrreversibleHarm}(x,s,a,s')]
-\lambda_{\mathrm{step}},
\end{aligned}
\label{eq:scorer}
\end{equation}
where $\Delta P_g(s,a,s')\triangleq P_g(s')-P_g(s)$ measures one-step progress toward subgoal $g$, $D_g(s')$ indicates subgoal completion, $\mathrm{Invalid}$ flags invalid actions, and $\mathrm{IrreversibleHarm}$ captures executor-level safety violations. The final term is a small per-step cost.

Each subgoal $g$ is parsed into a canonical type $\sigma=c(g)$ and instance slots $\zeta_g$ (e.g., object, receptacle, tool, product attribute, or option). We define normalized progress as
\begin{equation}
P_g(s)
=
\frac{\mathrm{rank}_{\sigma}(s;\zeta_g)}{n_{\sigma}}.
\label{eq:progress}
\end{equation}
For example, for $g=$ ``clean the cup'', $\sigma=\textsc{Clean}$ and $\zeta_g=\mathrm{cup}$:
\begin{equation}
\mathrm{rank}_{\textsc{Clean}}(s;\mathrm{cup})
=
\begin{cases}
0, & \text{cup not visible},\\
1, & \text{cup visible},\\
2, & \text{cup in inventory},\\
3, & \text{cup in inventory and agent at sink basin},\\
4, & \text{cup clean}.
\end{cases}
\label{eq:clean_rank}
\end{equation}
Analogous milestone functions are defined for other canonical subgoal types, such as \textsc{FindPick}, \textsc{Heat}, \textsc{Cool}, \textsc{Place}, \textsc{Search}, \textsc{Inspect}, and \textsc{Purchase}. Overall, this scorer provides dense partial credit for preparatory actions while remaining deterministic and inspectable.

\subsection{Anchored Branch Grouping for Low-Level Optimization} %\wei{这一节是主要创新，看起来有一些长，影响理解，建议可以加几个小标题}

The low-level executor should be optimized under a fixed active subgoal. A primitive action that is appropriate for one subgoal may be inappropriate for another, even under the same environment state. Therefore, HSPO constructs low-level comparison groups at the granularity of task--state--subgoal anchors.

ABG is applied after collecting a batch of complete task rollouts. Each rollout may contain multiple online subgoal segments; ABG scans all low-level decision points inside these complete trajectories.

For each transition $(s_t^{(i)},g_t^{(i)},a_t^{(i)},s_{t+1}^{(i)})$, ABG computes the anchor key
\begin{equation}
\kappa_t^{(i)}
=
(\kappa_x(x),\kappa_s(s_t^{(i)}),\kappa_g(g_t^{(i)})),
\label{eq:anchor}
\end{equation}
where $\kappa_x$ is a task signature, $\kappa_s$ is a canonical state signature, and $\kappa_g$ is a subgoal-instance signature. The subgoal instance is necessary because the same state and action may have different meanings under different active subgoals.

For each anchor $\kappa$, the natural group is
\begin{equation}
G^{\mathrm{nat}}(\kappa)
=
\{(s_i,a_i,s_i',\ell_i):\kappa_i=\kappa\},
\label{eq:natural_group}
\end{equation}
where $\ell_i$ is the old-policy log probability of the action span. Let $K(\kappa)=|G^{\mathrm{nat}}(\kappa)|$.

If the natural group is sufficiently large, it is used directly. If it is sparse, HSPO supplements it with one-step branched candidates. Let $K_{\min}$ be the minimum natural group size, $M_{\mathrm{target}}$ the target group size, and $B_{\max}$ the maximum branch budget per anchor. The number of branch candidates is
\begin{equation}
B(\kappa)
=
\begin{cases}
0, & K(\kappa)\geq K_{\min},\\
\min(M_{\mathrm{target}}-K(\kappa),B_{\max}), & K(\kappa)<K_{\min}.
\end{cases}
\label{eq:branch_budget}
\end{equation}
When $B(\kappa)>0$, HSPO restores a representative environment snapshot for anchor $\kappa$, samples $B(\kappa)$ actions from the old low-level policy under the same task, state, and active subgoal, and executes each for one environment step:


%\wei{representative 是指的恢复到哪一个snapshot，有很多。还有这里分支采样，不使用ht的话，是否会有问题，这里我是不确定的。}
%


\begin{equation}
a_{\mathrm{br}}^{(j)}
\sim
\pi_{\theta_{\mathrm{old}}}^{L}(\cdot\mid x,s_\kappa,g_\kappa),
\quad
s_\kappa \xrightarrow{a_{\mathrm{br}}^{(j)}} s_{\mathrm{br}}^{\prime(j)}.
\label{eq:branching}
\end{equation}
The final ABG group is
\begin{equation}
G^{\mathrm{ABG}}(\kappa)
=
G^{\mathrm{nat}}(\kappa)\cup G^{\mathrm{br}}(\kappa).
\label{eq:abg_group}
\end{equation}
Groups of size one are excluded from the relative low-level objective because no within-group comparison can be formed. These transitions receive no ABG update in the current batch, but they may still be covered by other anchors in later batches or regularized by the SFT loss.
%【已加】
 %\wei{这里加一句These transitions receive no low-level update from ABG, they may still be covered by other anchors or by the SFT loss. 解释那些transitions不参与更新可能引起的某些动作不被学习的问题，不会出现}


Each group element is scored by the milestone scorer:
\begin{equation}
q_i = S(x,g_\kappa,s_i,a_i,s_i').
\end{equation}
The group-normalized low-level advantage is
\begin{equation}
A_i^L
=
\frac{q_i-\mu(q_{G^{\mathrm{ABG}}(\kappa)})}
{\sigma(q_{G^{\mathrm{ABG}}(\kappa)})+\epsilon}.
\label{eq:low_adv}
\end{equation}
The low-level objective is a clipped policy-gradient loss applied only to action tokens:
\begin{equation}
\mathcal{L}^L(\theta)
=
-\mathbb{E}_{\kappa,i}
\left[
\min
\left(
\rho_i^L A_i^L,\ 
\mathrm{clip}(\rho_i^L,1-\epsilon,1+\epsilon)A_i^L
\right)
\right]
+
\beta_L\mathrm{KL}(\pi_\theta^L\|\pi_{\mathrm{ref}}^L),
\label{eq:low_loss}
\end{equation}
where
\begin{equation}
\rho_i^L
=
\frac{
\pi_\theta^L(a_i\mid x,s_i,g_\kappa)
}{
\pi_{\theta_{\mathrm{old}}}^L(a_i\mid x,s_i,g_\kappa)
}.
\end{equation}

\wei{$\pi_{\mathrm{ref}}^L$ denote the low-level policy after SFT warm-up phase}

Branched candidates are used only to construct training advantages. The real rollout trajectory always advances with the on-policy action originally sampled during interaction, not with the highest-scoring candidate. This avoids using the scorer as an inference-time action selector and prevents train-test mismatch.

\subsection{High-Level and Termination Objectives}

The planner is updated only at subgoal boundaries. At boundary $b_k$, the model emits subgoal $g_k$, and the executor acts under this subgoal until the next boundary. The macro transition is $(s_{b_k},g_k,s_{b_{k+1}})$. We define the high-level reward as 

\begin{equation}
R_k^H
=
\mathbf{1}[k=K-1]R_{\mathrm{task}}
-
\beta C_{\mathrm{side}}(x,s_{b_k},s_{b_{k+1}})
-
\eta C_{\mathrm{red}}(g_k,h_{b_k}),
\label{eq:high_reward}
\end{equation}
where $C_{\mathrm{side}}$ penalizes globally harmful state changes during the segment and $C_{\mathrm{red}}$ penalizes redundant subgoals. A macro critic $V^H$ estimates boundary-state values:
\begin{equation}
\delta_k^H
=
R_k^H+\gamma_H V^H(s_{b_{k+1}})-V^H(s_{b_k}),
\quad
A_k^H
=
\sum_{j=k}^{K-1}
(\gamma_H\lambda_H)^{j-k}\delta_j^H.
\label{eq:macro_gae}
\end{equation}
The planner loss is
\begin{equation}
\mathcal{L}^H(\theta)
=
-\mathbb{E}_k
\left[
\min
\left(
\rho_k^H A_k^H,\ 
\mathrm{clip}(\rho_k^H,1-\epsilon,1+\epsilon)A_k^H
\right)
\right]
+
\beta_H\mathrm{KL}(\pi_\theta^H\|\pi_{\mathrm{ref}}^H),
\label{eq:high_loss}
\end{equation}
where
\begin{equation}
\rho_k^H
=
\frac{
\pi_\theta^H(g_k\mid x,s_{b_k},h_{b_k})
}{
\pi_{\theta_{\mathrm{old}}}^H(g_k\mid x,s_{b_k},h_{b_k})
}.
\end{equation}
This loss updates only subgoal tokens.

%\wei{这地方两个loss之间加一点过渡性质的描述，Unlike the high‑level planner, which is updated via macro PPO over segment boundaries, the termination decision (the <switch> span) is trained with supervised lagged targets. This is because termination depends only on local subgoal completion, not on long‑horizon returns, and supervised learning is more stable when the executor is still unreliable.}


The switch span is trained with lagged subgoal-completion supervision. If action $a_t$ completes the current subgoal, the next model call should switch; the completion action itself still belongs to the current segment. Thus,
\begin{equation}
y_{t+1}^T
=
\begin{cases}
\mathsf{SWITCH}, & D_{g_t}(s_{t+1})=1,\\
\mathsf{KEEP}, & D_{g_t}(s_{t+1})=0.
\end{cases}
\label{eq:switch_target}
\end{equation}
The termination loss is
\begin{equation}
\mathcal{L}^T(\theta)
=
-\mathbb{E}_t
\log
\pi_\theta^T(y_t^T\mid x,s_t,g_{t-1},h_t),
\label{eq:switch_loss}
\end{equation}
applied only to switch tokens.

\subsection{Token-Level Credit Masks and Total Objective}
\label{sec:token_masks}

HSPO is implemented with a single autoregressive LLM, but different parts of its structured response correspond to different decision levels. To prevent credit from one level from directly updating another, we apply token-level credit masks to the output spans. Let
\begin{equation}
m_i^T=\mathbf{1}[y_i\in \texttt{<switch>}],\quad
m_i^H=\mathbf{1}[y_i\in \texttt{<subgoal>}],\quad
m_i^L=\mathbf{1}[y_i\in \texttt{<action>}],
\label{eq:token_masks}
\end{equation}
where $m_i^T$, $m_i^H$, and $m_i^L$ indicate whether token $y_i$ belongs to the termination, subgoal, or action span, respectively. Each objective is computed only over its corresponding span: high-level macro advantages update subgoal tokens, low-level ABG advantages update action tokens, and termination supervision updates switch tokens.

Equivalently, for a response $y$, the masked log probabilities are
\begin{equation}
\begin{aligned}
\log \pi_\theta^T(y) &= \sum_i m_i^T \log \pi_\theta(y_i\mid y_{<i},c), \\
\log \pi_\theta^H(y) &= \sum_i m_i^H \log \pi_\theta(y_i\mid y_{<i},c), \\
\log \pi_\theta^L(y) &= \sum_i m_i^L \log \pi_\theta(y_i\mid y_{<i},c).
\end{aligned}
\label{eq:masked_logprob}
\end{equation}
These masked log probabilities are used in the termination, high-level, and low-level losses introduced above. The masking does not create separate models; rather, it routes different learning signals to different output spans of the same model.

However, token masks alone do not completely isolate the underlying parameters, since all spans are generated by a shared LLM. In particular, low-level action updates may still indirectly change the distribution of future subgoal tokens through shared representations. To reduce such planner drift during executor optimization, we regularize the subgoal-token distribution against the supervised Plan-Execute reference policy $\pi_{\mathrm{SFT}}$: %\wei{这里的$\pi_{\mathrm{SFT}}$应该就是前面没定义的$\pi_{\mathrm{ref}}^L$吧，前面我加了定义，如果是一个意思的话，就统一一下符号}
\begin{equation}
\mathcal{L}_{\mathrm{sgKL}}(\theta)
=
\mathbb{E}
\left[
\sum_i m_i^H
\mathrm{KL}
\left(
\pi_\theta(\cdot\mid y_{<i},c)
\|
\pi_{\mathrm{SFT}}(\cdot\mid y_{<i},c)
\right)
\right].
\label{eq:subgoal_kl}
\end{equation}
This term stabilizes subgoal generation while the executor is being optimized. It should be interpreted as distributional stabilization rather than strict parameter-level freezing.

The final training objective combines high-level planning, low-level execution, termination learning, value learning, supervised regularization, and subgoal-distribution stabilization:
\begin{equation}
\mathcal{L}_{\mathrm{HSPO}}
=
\mathcal{L}^{H}
+
\alpha_L \mathcal{L}^{L}
+
\alpha_T \mathcal{L}^{T}
+
\alpha_V \mathcal{L}_{V^H}
+
\alpha_{\mathrm{SFT}}\mathcal{L}_{\mathrm{SFT}}
+
\beta_{\mathrm{sg}}\mathcal{L}_{\mathrm{sgKL}}.
\label{eq:total_loss}
\end{equation}
Here, $\mathcal{L}^{H}$ optimizes subgoal selection, $\mathcal{L}^{L}$ optimizes subgoal-conditioned action execution, $\mathcal{L}^{T}$ trains online subgoal termination, and $\mathcal{L}_{V^H}$ trains the macro critic. The SFT term preserves the structured Plan-Execute format, while the subgoal KL term reduces unintended drift of the planner during low-level updates.

This objective realizes the desired credit separation. A subgoal can receive negative high-level advantage because it is globally harmful, while an action inside that segment can receive positive low-level advantage because it correctly advances the active subgoal. Since the two learning signals are applied to different output spans, HSPO can suppress harmful subgoal choices without erasing reusable executor skills.

HSPO is trained with a staged optimization procedure. We first perform a lightweight Plan-Execute supervised warm-up to initialize the structured response format and basic execution behavior. We then construct and validate the state-grounded milestone scorer before RL. During low-level training, complete task rollouts are collected, ABG groups are built from task--state--subgoal anchors, sparse groups are supplemented with one-step branches, and action tokens are updated with the low-level objective; switch tokens are trained with lagged completion supervision, while subgoal-distribution regularization reduces unintended planner drift. Finally, after the executor and termination behavior are stabilized, the planner is optimized over macro transitions with the high-level objective and macro critic. Algorithm~\ref{alg:hspo} in Appendix~\ref{app:hspo_algorithm} summarizes the full staged training procedure.
%\wei{这一段略长，可以根据最后篇幅长度调整，凝练一些}


%伪代码放附录2
\section{Experiments}
\label{sec:experiments}

We evaluate HSPO from two complementary perspectives. First, we examine whether decoupling high-level subgoal credit from low-level subgoal-conditioned execution credit improves long-horizon task performance compared with flat RL, group-based RL, and prompt-based agent baselines. Second, we directly test the proposed failure mode, \emph{intra-subgoal credit collapse}, by measuring whether the executor preserves locally correct skills even when the corresponding high-level subgoals are globally suboptimal. Accordingly, our evaluation includes standard task-level metrics as well as targeted diagnostics for subgoal execution, milestone-scorer quality, anchor validity, planner drift, and training cost.

\subsection{Experimental Setup}
\label{sec:exp_setup}

\paragraph{Benchmarks.}
We evaluate on two long-horizon interactive agent benchmarks. \textbf{ALFWorld}~\cite{shridhar2021alfworld} is a text-based embodied household environment that requires agents to perform multi-step tasks through language actions. It contains six task categories: Pick \& Place, Examine in Light, Clean \& Place, Heat \& Place, Cool \& Place, and Pick Two \& Place. We evaluate on 134 unseen test instances. \textbf{WebShop}~\cite{yao2022webshop} is a simulated e-commerce environment where agents must search for products, inspect item pages, match user-specified attributes, select options, and issue a purchase action. We evaluate on 500 test instructions with a maximum episode length of 15 turns.

\paragraph{Models.}
We evaluate HSPO and all trainable RL baselines with two base model sizes: \textbf{Qwen2.5-1.5B-Instruct} and \textbf{Qwen2.5-7B-Instruct}. For each model size, all trainable methods are initialized from the same Plan-Execute SFT checkpoint, use the same training data, and follow the same rollout and evaluation protocol. Prompt-based baselines are evaluated separately as reference systems and do not use our SFT or RL training data.

\paragraph{Baselines.}
We compare HSPO with both prompt-based and trainable baselines. The prompt-based baselines include \textbf{ReAct}~\cite{yao2022react}, which interleaves reasoning traces and environment actions, and \textbf{Reflexion}~\cite{shinn2023reflexion}, which augments ReAct-style interaction with verbal self-reflection. Unless otherwise specified, these prompt-based baselines are evaluated with GPT-4o under the same environment protocol.

The trainable baselines include \textbf{SFT only}, which uses the Plan-Execute warm-up checkpoint without RL; \textbf{PPO}~\cite{schulman2017proximal} with a learned value critic; \textbf{RLOO}~\cite{ahmadian2024back}; trajectory-level \textbf{GRPO}~\cite{shao2024deepseekmath}; and \textbf{GiGPO}~\cite{NEURIPS2025_420c9f77}, which constructs grouped advantages from repeated environment states. These baselines cover supervised warm-up, critic-based policy optimization, critic-free trajectory-level group optimization, and state-anchored step-level group optimization. Full implementation details, including prompt templates and baseline-specific training settings, are provided in Appendix~\ref{app:implementation_details}.

\paragraph{Training and implementation details.}
Unless otherwise specified, all trainable methods use the same actor model, rollout budget, maximum episode length, and evaluation protocol within each model size. The actor learning rate is $1{\times}10^{-6}$, and the critic learning rate, where applicable, is $1{\times}10^{-5}$. The maximum prompt length is 4096 tokens and the maximum response length is 512 tokens. The maximum episode length is 50 steps for ALFWorld and 15 turns for WebShop. For group-based methods, the trajectory group size is 8. For HSPO, the target ABG group size is $M_{\mathrm{target}}=4$, and the maximum branch budget is $B_{\max}=4$ unless otherwise stated. The mini-batch size is 256, the rollout temperature is 1.0, and the evaluation temperature is 0.4.

For HSPO, we train the low-level executor and termination module for 150 RL updates, followed by 100 macro-level planner updates. The PPO clipping coefficient is $\epsilon_{\mathrm{clip}}=0.2$, and the advantage-normalization constant is $\epsilon_{\mathrm{adv}}=10^{-6}$. The loss coefficients are $\alpha_L=1.0$, $\alpha_T=0.5$, $\alpha_V=0.5$, $\alpha_{\mathrm{SFT}}=0.1$, and $\beta_L=\beta_H=0.01$. The subgoal-distribution regularization coefficient $\beta_{\mathrm{sg}}$ is selected on the validation split. All experiments are conducted on \texttt{[number]} \texttt{[GPU type]} GPUs. Additional hardware details, wall-clock training time, and memory usage are reported in Appendix~\ref{app:compute_cost}.

\paragraph{Metrics.}
For ALFWorld, we report task success rate, defined as the fraction of test tasks completed within 50 environment steps. We also report category-wise success rates for the six task types. When subgoal annotations are available, we additionally report progress rate using AgentBoard-style subgoal progress annotations~\cite{ma2024agentboard}. For WebShop, we report the official task score on a 0--100 scale and exact success rate, where success means achieving a score of 100. All trainable methods are evaluated over three random seeds. We report mean performance and standard deviation; per-seed results are included in Appendix~\ref{app:per_seed_results} to make the stability of each method transparent.


\begin{table*}[t]
\centering
\small
\setlength{\tabcolsep}{3.2pt}
\renewcommand{\arraystretch}{1.08}
\resizebox{\textwidth}{!}{
\begin{tabular}{lll|ccccccc|cc}
\toprule
\multicolumn{3}{c|}{\textbf{Setting}} 
& \multicolumn{7}{c|}{\textbf{ALFWorld}} 
& \multicolumn{2}{c}{\textbf{WebShop}} \\
\cmidrule(lr){1-3} \cmidrule(lr){4-10} \cmidrule(l){11-12}
\textbf{Type} & \textbf{Model} & \textbf{Method} 
& \textbf{Pick} & \textbf{Look} & \textbf{Clean} & \textbf{Heat} & \textbf{Cool} & \textbf{Pick2} & \textbf{All} 
& \textbf{Score} & \textbf{Succ.} \\
\midrule
PE & GPT-4o & ReAct~\cite{yao2022react}      
    &  &  &  &  &  &  &  &  &  \\
PE & GPT-4o & Reflexion~\cite{shinn2023reflexion}  
    &  &  &  &  &  &  &  &  &  \\
\midrule
SFT & Qwen2.5-1.5B & SFT only
    &  &  &  &  &  &  &  &  &  \\
RL & Qwen2.5-1.5B & PPO~\cite{schulman2017proximal}        
    &  &  &  &  &  &  &  &  &  \\
RL & Qwen2.5-1.5B & RLOO~\cite{ahmadian2024back}       
    &  &  &  &  &  &  &  &  &  \\
RL & Qwen2.5-1.5B & GRPO~\cite{shao2024deepseekmath}       
    &  &  &  &  &  &  &  &  &  \\
RL & Qwen2.5-1.5B & GiGPO~\cite{NEURIPS2025_420c9f77}      
    &  &  &  &  &  &  &  &  &  \\
RL & Qwen2.5-1.5B & \textbf{HSPO (Ours)} 
    &  &  &  &  &  &  &  &  &  \\
\midrule
SFT & Qwen2.5-7B & SFT only
    &  &  &  &  &  &  &  &  &  \\
RL & Qwen2.5-7B & PPO~\cite{schulman2017proximal}        
    &  &  &  &  &  &  &  &  &  \\
RL & Qwen2.5-7B & RLOO~\cite{ahmadian2024back}       
    &  &  &  &  &  &  &  &  &  \\
RL & Qwen2.5-7B & GRPO~\cite{shao2024deepseekmath}       
    &  &  &  &  &  &  &  &  &  \\
RL & Qwen2.5-7B & GiGPO~\cite{NEURIPS2025_420c9f77}      
    &  &  &  &  &  &  &  &  &  \\
RL & Qwen2.5-7B & \textbf{HSPO (Ours)} 
    &  &  &  &  &  &  &  &  &  \\
\bottomrule
\end{tabular}
}
\caption{\textbf{Main results on ALFWorld and WebShop.}
For ALFWorld, we report success rates by task category and overall average. For WebShop, we report the official task score and exact success rate. Trainable methods are evaluated with both Qwen2.5-1.5B-Instruct and Qwen2.5-7B-Instruct and averaged over three random seeds. Standard deviations and per-seed results are reported in Appendix~\ref{app:per_seed_results}.}
\label{tab:main_results}
\end{table*}


\begin{figure*}[t]
\centering
\includegraphics[width=\textwidth]{figures/fig3-traj.png}
\caption{
\textbf{Diagnostic analysis on ALFWorld demonstration trajectories.}
(a) Trajectory length distribution across task types. 
(b) Number of induced subgoal segments per trajectory, showing that multi-stage tasks require multiple online subgoal decisions.
(c) Subgoal-type composition by task type, computed over subgoal segments rather than raw low-level steps.
(d) Intrinsic sanity check of the deterministic milestone scorer. For each subgoal type, bars report the mean scorer output over real transitions that increase the milestone rank ($\Delta\mathrm{rank}>0$) or do not increase it ($\Delta\mathrm{rank}\leq0$).
}

\label{fig:hspo_diagnostics}
\end{figure*}

\subsection{Main Results}
\label{sec:main_results}

Table~\ref{tab:main_results} reports the main comparison on ALFWorld and WebShop across two model scales. HSPO is designed to improve long-horizon agent learning by assigning high-level and low-level credit to different decision spans. Therefore, we expect the largest gains on tasks where successful completion requires both choosing an appropriate sequence of subgoals and executing each subgoal through multiple environment actions, such as Clean \& Place, Heat \& Place, Cool \& Place, and WebShop product selection.

Compared with flat RL baselines such as PPO, RLOO, and GRPO, HSPO explicitly separates planner optimization from executor optimization. This prevents final task outcomes from being uniformly propagated to all low-level actions. Compared with GiGPO, HSPO further conditions low-level comparison groups on the active subgoal, so that actions are compared under the same local intent rather than merely under the same environment state. The results in Table~\ref{tab:main_results} are used to test whether this additional subgoal conditioning translates into stronger task-level performance.

The two model scales further allow us to examine whether HSPO's benefits come only from increased model capacity or from the proposed credit-assignment mechanism itself. If HSPO consistently improves over the corresponding baselines within both Qwen2.5-1.5B and Qwen2.5-7B settings, it would suggest that subgoal-conditioned credit separation provides benefits beyond simple scaling. In addition to overall performance, category-wise ALFWorld results help diagnose where HSPO provides the most benefit. Improvements on manipulation-heavy categories indicate stronger low-level execution, while improvements on categories requiring object transformation and placement suggest better coordination between high-level subgoal planning and low-level action execution.

\subsection{Ablation Studies}
\label{sec:ablations}

We conduct ablation studies to isolate the contribution of each major component in HSPO. Following the main experimental protocol, all ablations are performed on ALFWorld with Qwen2.5-1.5B-Instruct, using the same SFT initialization, rollout budget, and training steps. We report task success rate and progress rate, and include executor skill retention to directly measure whether locally correct low-level behaviors are preserved under globally suboptimal subgoals.

\begin{table}[t]
\centering
\small
\setlength{\tabcolsep}{5pt}
\renewcommand{\arraystretch}{1.08}
\begin{tabular}{l|ccc}
\toprule
\textbf{Method}
& \textbf{Success (\%)}
& \textbf{Progress (\%)}
& \textbf{Skill Retention (\%)} \\
\midrule
HSPO
&  &  &  \\
\quad w/o ABG
&  &  &  \\
\quad w/o subgoal anchor
&  &  &  \\
\quad w/o branching
&  &  &  \\
\quad w/o milestone scorer
&  &  &  \\
\quad w/o token routing
&  &  &  \\
\quad w/o subgoal KL
&  &  &  \\
\quad w/o macro planner optimization
&  &  &  \\
\bottomrule
\end{tabular}
\caption{\textbf{Ablation results on ALFWorld.}
All variants use Qwen2.5-1.5B-Instruct. Success measures final task completion, Progress measures annotated subgoal progress, and Skill Retention measures whether the executor preserves locally correct actions under globally suboptimal subgoals.}
\label{tab:ablation}
\end{table}

\paragraph{Ablation variants.}
\textbf{w/o ABG} removes Anchored Branch Grouping and optimizes low-level actions with segment-level feedback, testing whether subgoal-conditioned local comparison is necessary.
\textbf{w/o subgoal anchor} groups actions only by task and state, removing $\kappa_g(g_t)$ from the anchor. This variant tests whether state-level grouping is sufficient or whether the active subgoal must be included.
\textbf{w/o branching} uses only naturally repeated anchors and disables one-step branch supplementation, testing whether branching improves the coverage of sparse low-level comparison groups.
\textbf{w/o milestone scorer} replaces the state-grounded process reward with outcome-based segment feedback, testing whether one-step progress toward the active subgoal is needed for low-level credit assignment.
\textbf{w/o token routing} applies all losses to the full response instead of their corresponding \texttt{<switch>}, \texttt{<subgoal>}, and \texttt{<action>} spans, testing whether token-level credit separation is necessary.
\textbf{w/o subgoal KL} removes $\mathcal{L}_{\mathrm{sgKL}}$, testing whether planner drift occurs during executor optimization.
\textbf{w/o macro planner optimization} removes the high-level objective $\mathcal{L}^H$, testing whether improving the executor alone is sufficient for long-horizon task success.

\paragraph{Discussion.}
These ablations are designed to test the main claims of HSPO. If removing ABG or the milestone scorer reduces skill retention, it indicates that low-level actions again inherit global or segment-level credit, leading to intra-subgoal credit collapse. If removing the subgoal anchor hurts performance, it suggests that actions should be compared under the same task--state--subgoal intent rather than under state alone. If removing token routing or subgoal KL degrades performance, it shows that separating credit at the token level and stabilizing the planner are important when a single LLM shares parameters across planning and execution.

\subsection{Analysis}


\section{Limitations}
\wei{nips鼓励有这个章节，这里可以大方承认一些短板，也可以提前封住审稿人的口}

HSPO is designed for \textbf{milestone-observable interactive environments}: environments where intermediate progress toward subgoals can be extracted from structured or semi-structured state and where state can be cloned or replayed for one-step branching.

Full ABG requires environment snapshots for under-populated anchors when $B_{\max} > 0$. This is practical in symbolic text games and simulators with serializable state. It is more difficult in session-based web environments, where session snapshots must include page state, product state, selected options, cart state, and backend randomness. In irreversible real-world environments, external APIs, production browsers, or OS-level agents, full branching may be infeasible; the zero-branch configuration $B_{\max} = 0$ is the appropriate fallback.

The scorer requires per-environment engineering. Structured annotations can reduce design cost, but parser, instance extractor, state extractor, rank functions, invalid checker, and safety predicates must still be implemented. HSPO does not claim zero-cost transfer of the scorer across environments.

The canonical subgoal vocabulary $\Sigma$ is currently predefined. Open-vocabulary subgoals would require a robust subgoal-type classifier or learned parser, which may reintroduce uncertainty.

Finally, token masks and KL regularization stabilize span-level outputs but do not fully isolate parameters in a shared LLM. Span-specific adapters or parameter freezing may provide stronger isolation but introduce additional engineering and capacity trade-offs.
%\wei{当前版本和paperreview.ai交互的审稿人提出的问题，这个审稿比较严苛，我先放到这，你先不用看 1. Please provide the complete main results (success/score with error bars over seeds) and ablations referenced in Sections 4.2–4.4. Which components of HSPO (ABG, branching, token masks, sg-KL, high-level macro PPO) are most critical? 2.How are κ-state signatures and subgoal-instance signatures constructed in practice for ALFWorld and WebShop? Are they rule-based, learned, or based on environment metadata? Can you report collision rates and anchor purity? 3.How do you detect Dg (subgoal completion) and compute Cside and Cred? Are these derived solely from the milestone definitions or additional environment-specific oracles? Please clarify any external annotations used. 4. What is the empirical correlation between the milestone scorer S and final task return across states/subgoals? Do actions that maximize S consistently lead to better end-task success, or do you observe scorer hacking? 5.How feasible is environment snapshot restoration and one-step branching on platforms without determinism or exact-state save/restore (e.g., real robots, some web environments)? Can HSPO operate without branching, and with what performance penalty? 6. Please include sensitivity analyses for Mtarget, Bmax, Kmin, and scorer weights (λinv, λhard, ηdone), as well as for the sg-KL coefficient. How robust is HSPO to these hyperparameters? 7.Can you compare HSPO to HiMAC under identical settings (same base LLM, data splits, and compute) and discuss trade-offs between online switching (HSPO) and blueprint-based alternating training (HiMAC)? 8.How much cross-span drift remains despite token masks and sg-KL? Can you report quantitative measures of planner distribution shift during low-level training, with/without sg-KL?}




\bibliographystyle{unsrt}
\bibliography{Ref}

\appendix %\wei{附录补充prompt格式举例、算法、如果有信用坍塌的验证或者一定的理论证明更好、还有较完整的超参数表}

\section{Full Milestone Rank Definitions}

\subsection{ALFWorld}

\textbf{FindPick.} $\mathrm{rank}(s) \in \{0, 1, 2\}$: 0 if the target object's containing receptacle has not been opened or examined; 1 if the target object is visible (i.e., its location is observable); 2 if the target object is in the agent's inventory.

\textbf{Clean.} Defined in Eq.~(5).

\textbf{Heat.} $\mathrm{rank}(s) \in \{0, 1, 2, 3, 4\}$: 0 target not visible, 1 target visible, 2 target in inventory, 3 target in inventory and agent at microwave, 4 target heated.

\textbf{Cool.} $\mathrm{rank}(s) \in \{0, 1, 2, 3, 4\}$: analogous to Heat but with fridge and cool state.

\textbf{Place.} $\mathrm{rank}(s) \in \{0, 1, 2, 3\}$: 0 target not held, 1 target in inventory, 2 target in inventory and target receptacle visible, 3 target placed in target receptacle.

\textbf{Examine.} $\mathrm{rank}(s) \in \{0, 1, 2, 3\}$: 0 target not visible, 1 target visible, 2 target in inventory and light source visible, 3 target examined in light (light on, target held in the light source's vicinity).

\subsection{WebShop}

\textbf{Search.} $\mathrm{rank}(s) \in \{0, 1, 2\}$: 0 no search performed, 1 search performed but query coverage of required attributes $< 50\%$, 2 search performed with $\geq 50\%$ coverage.

\textbf{Inspect.} $\mathrm{rank}(s) \in \{0, 1, 2\}$: 0 no product clicked, 1 product clicked, 2 product clicked and detail page reached.

\textbf{SelectOption.} $\mathrm{rank}(s) \in \{0, 1, 2, 3\}$: 0 no options selected, 1 some options selected, 2 all required options selected, 3 all options match goal attributes.

\textbf{Purchase.} $\mathrm{rank}(s) \in \{0, 1, 2\}$: 0 not at purchase page, 1 at purchase page, 2 purchase confirmed.
\clearpage

\section{HSPO Training Algorithm}
\label{app:hspo_algorithm}

\begin{algorithm}[H]
\small
\caption{HSPO with Anchored Branch Grouping}
\label{alg:hspo}
\begin{algorithmic}[1]
\Require Policy $\pi_\theta$, SFT reference $\pi_{\mathrm{SFT}}$, macro critic $V^H$, scorer $S$, environment $\mathrm{Env}$, ABG parameters $K_{\min},M_{\mathrm{target}},B_{\max}$
\Ensure Trained policy $\pi_\theta$
\State \textbf{Phase I:} Train $\pi_\theta$ on structured Plan-Execute demonstrations with $\mathcal{L}_{\mathrm{SFT}}$.
\State \textbf{Phase II:} Define canonical subgoal types, parsers, milestone ranks, invalid checks, and safety checks.
\State Validate the scorer on held-out demonstration transitions.
\State \textbf{Phase III: Low-level ABG optimization}
\For{each training batch}
    \State Roll out complete task trajectories with $\pi_\theta$.
    \State Store all transitions $(x,s_t,g_t,a_t,s_{t+1})$, old log probabilities, and snapshots.
    \State Build anchor table using $\kappa_t=(\kappa_x(x),\kappa_s(s_t),\kappa_g(g_t))$.
    \For{each anchor $\kappa$}
        \State Construct $G^{\mathrm{nat}}(\kappa)$.
        \If{$|G^{\mathrm{nat}}(\kappa)| < K_{\min}$}
            \State Branch up to $B_{\max}$ one-step candidate actions toward $M_{\mathrm{target}}$.
        \EndIf
        \State Form $G^{\mathrm{ABG}}(\kappa)$ and compute scorer values $q_i$.
        \State Normalize group scores to obtain $A_i^L$.
    \EndFor
    \State Compute lagged switch targets from real on-policy transitions.
    \State Update action spans with $\mathcal{L}^L$, switch spans with $\mathcal{L}^T$, and apply $\mathcal{L}_{\mathrm{sgKL}}$.
\EndFor
\State \textbf{Phase IV: High-level planner optimization}
\For{each training batch}
    \State Roll out complete trajectories and extract macro transitions $(s_{b_k},g_k,s_{b_{k+1}})$.
    \State Compute high-level rewards $R_k^H$ and macro advantages $A_k^H$.
    \State Update subgoal spans with $\mathcal{L}^H$ and train $V^H$.
\EndFor
\State \Return $\pi_\theta$
\end{algorithmic}
\end{algorithm}

\section{Hyperparameters and Implementation Details}

All experiments are conducted on Ubuntu Linux (Ubuntu 22.04 LTS) with NVIDIA A800-80GB GPUs, using a single-GPU setting by default with an FSDP + vLLM hybrid training/inference pipeline. We use Qwen2.5-1.5B-Instruct as the backbone model and set random seed to 0.

For ALFWorld, the main settings are: train/validation batch size $8/16$, group size $4$, maximum episode length $50$, prompt/response length $2048/512$, actor learning rate $1\times 10^{-6}$, and total training epochs $200$. For WebShop, we use train/validation batch size $4/16$, group size $2$, maximum episode length $15$, prompt/response length $4096/512$, actor learning rate $1\times 10^{-6}$, and total training epochs $200$.

Across both environments, we use gradient clipping with threshold $1.0$, constant warmup-style learning-rate scheduling, and HSPO defaults for hierarchical optimization (e.g., low/high PPO clip $\epsilon=0.2$, $(\gamma_H,\lambda_H)=(0.99,0.95)$, and ABG grouping with min/target/max branch $=2/4/4$). Checkpoints are saved every 25 epochs and validation is run every 5 epochs.


\section{Prompt Templates}
\label{app:prompts}

We use ReAct-style prompts for flat baselines and Plan-Execute prompts for HSPO. In HSPO, each step must return exactly three tags:
\texttt{<switch>} (\texttt{KEEP}/\texttt{SWITCH}),
\texttt{<subgoal>}, and
\texttt{<action>}.
Below we provide representative templates (slightly shortened for readability, while preserving the exact output interface constraints).

\subsection{ALFWorld Prompts}
\label{app:alfworld_prompts}

\paragraph{Flat baseline (ReAct-style).}
\begin{lstlisting}
You are an expert agent in ALFWorld.
Task: {task_description}
Current observation: {current_observation}
Admissible actions: [{admissible_actions}]

Think step-by-step, then output exactly one action:
<action>ONE admissible action</action>
\end{lstlisting}

\paragraph{HSPO (Plan-Execute).}
\begin{lstlisting}
You are an expert agent operating in the ALFRED Embodied Environment.
Your overall task is: {task_description}
You have already taken {step_count} step(s).
Most recent {history_length} observations and actions:
{action_history}

Your current observation is:
{current_observation}

Your current sub-goal is:
{current_subgoal}

Your admissible actions are:
[{admissible_actions}]

At EVERY step, output EXACTLY THREE blocks in order:
<switch>KEEP or SWITCH</switch>
<subgoal>the sub-goal you will follow next</subgoal>
<action>EXACTLY one ADMISSIBLE action</action>
\end{lstlisting}

\subsection{WebShop Prompts}
\label{app:webshop_prompts}

\paragraph{Flat baseline (ReAct-style).}
\begin{lstlisting}
You are an autonomous shopping assistant in WebShop.
Task: {task_description}
Current observation: {current_observation}
Available actions:
[
{available_actions}
]

Reason briefly, then output one action:
<action>ONE valid action (search[...] or click[...])</action>
\end{lstlisting}

\paragraph{HSPO (Plan-Execute).}
\begin{lstlisting}
You are an expert autonomous agent operating in WebShop.
Your task is to: {task_description}
You have already taken {step_count} step(s).
Recent {history_length} observations/actions: {action_history}
Current observation: {current_observation}
Current sub-goal: {current_subgoal}
Admissible actions:
[
{available_actions}
]

Output EXACTLY THREE blocks, no extra text:
<switch>KEEP or SWITCH</switch>
<subgoal>the sub-goal you will follow next</subgoal>
<action>EXACTLY one ADMISSIBLE action</action>
\end{lstlisting}

\section{Qualitative Trajectories}
\label{app:qualitative}

We show representative successful trajectories from the released setup to illustrate how HSPO separates subgoal switching (\texttt{<switch>}/\texttt{<subgoal>}) from execution (\texttt{<action>}).

\subsection{ALFWorld Example}
\label{app:alfworld_trajectory}

\begin{lstlisting}
Task: put a clean mug in the microwave.

Step 1
Obs: You are in the kitchen. You see a sink, a microwave, and a mug on the counter.
Output:
<switch>SWITCH</switch>
<subgoal>Pick up the mug and move to sink for cleaning</subgoal>
<action>take mug from countertop</action>

Step 2
Obs: You hold the mug. Sink is visible.
Output:
<switch>KEEP</switch>
<subgoal>Pick up the mug and move to sink for cleaning</subgoal>
<action>go to sink</action>

Step 3
Obs: At sink. Mug in inventory.
Output:
<switch>SWITCH</switch>
<subgoal>Clean the mug at the sink</subgoal>
<action>clean mug with sink</action>

Step 4
Obs: Mug is clean. Microwave is visible.
Output:
<switch>SWITCH</switch>
<subgoal>Place the clean mug into microwave</subgoal>
<action>open microwave</action>

Step 5
Obs: Microwave open.
Output:
<switch>KEEP</switch>
<subgoal>Place the clean mug into microwave</subgoal>
<action>put mug in microwave</action>
\end{lstlisting}

\subsection{WebShop Example}
\label{app:webshop_trajectory}

\begin{lstlisting}
Task: find a men's shirt (navy, size M, price < $50) and buy it.

Step 1
Obs: Search page.
Output:
<switch>SWITCH</switch>
<subgoal>Search with key attributes to narrow candidate items</subgoal>
<action>search[men's shirt navy medium price < 50]</action>

Step 2
Obs: Result list (multiple items).
Output:
<switch>SWITCH</switch>
<subgoal>Open a likely matching product and inspect details</subgoal>
<action>click[B07XDRVVYM]</action>

Step 3
Obs: Product page with options color/size.
Output:
<switch>SWITCH</switch>
<subgoal>Select required color and size options</subgoal>
<action>click[navy]</action>

Step 4
Obs: Color selected; size options visible.
Output:
<switch>KEEP</switch>
<subgoal>Select required color and size options</subgoal>
<action>click[medium]</action>

Step 5
Obs: Required options selected; Buy Now available.
Output:
<switch>SWITCH</switch>
<subgoal>Complete purchase on the matched product</subgoal>
<action>click[buy now]</action>
\end{lstlisting}




\newpage
\input{checklist.tex}

\end{document}
